\documentclass[10pt,twocolumn]{article}

\usepackage{cvpr}
\usepackage{times}
\usepackage{epsfig}
\usepackage{graphicx}
\usepackage{amsmath}
\usepackage{amssymb}
\usepackage[pagebackref,breaklinks,colorlinks]{hyperref}

\usepackage[utf8]{inputenc}
\usepackage[T1]{fontenc}
\usepackage{booktabs}
\usepackage{amsfonts}
\usepackage{mathtools}
\usepackage{algorithm}
\usepackage{algorithmic}
\usepackage{subcaption}

\DeclareMathOperator*{\argmin}{arg\,min}
\DeclareMathOperator*{\argmax}{arg\,max}

\newcommand{\R}{\mathbb{R}}
\newcommand{\E}{\mathbb{E}}
\newcommand{\Var}{\mathrm{Var}}
\newcommand{\Cov}{\mathrm{Cov}}
\newcommand{\KL}{D_{\mathrm{KL}}}

\cvprfinalcopy

\title{Paper Title Goes Here}

\author{
  Author One \\
  Affiliation \\
  \texttt{author1@example.com} \\
  \And
  Author Two \\
  Affiliation \\
  \texttt{author2@example.com}
}

\begin{document}

\maketitle

\begin{abstract}
Abstract goes here.
\end{abstract}

\section{Introduction}
Introduction goes here.

\section{Related Work}
Related work goes here.

\section{Method}
Method goes here.

\section{Experiments}
Experiments go here.

\section{Conclusion}
Conclusion goes here.

\bibliographystyle{plain}
\bibliography{references}

\end{document}
